\textbf{ছবিটি আসলে কোথায় থাকে?} এই পাঠের সাধারণ উদাহরণে \textbf{src} ছবির ফাইলের ঠিকানা দেয়; HTML ফাইলে ছবির সব বাইট লিখে রাখা হয় না। ব্রাউজার পথটি বর্তমান পেজের অবস্থানের সঙ্গে মিলিয়ে ফাইল খুঁজে লোড করে। তাই HTML পাঠানোর পাশাপাশি সংশ্লিষ্ট ছবিও সঠিক ফোল্ডারে রাখতে হবে। \textbf{alt} ছবির কাজ বা প্রয়োজনীয় তথ্য বোঝাবে; নিছক সাজসজ্জার ছবিতে \texttt{alt=""} উপযুক্ত হতে পারে।
